\section{Discussion}
\label{sec:discussion}

\begin{quote}\small\itshape
\begin{itemize}
\item Empty for the moment
\end{itemize}
\end{quote}

\doubt{Nothing drafted as prose here yet, per your note. Once
Section~\ref{sec:results} has real numbers, this section should close the
loop back to Section~\ref{sec:introduction}'s motivation. Suggested threads
to pick up:}

\begin{itemize}
\item How much smaller a vocabulary does structured (greedy / degree /
centrality) selection buy over the "keep everything" or "random" baselines
from the Introduction's gap paragraph, at what cost in semantic coverage or
distance score -- i.e.\ turn Section~\ref{sec:results:trend}/\ref{sec:results:families}
into a direct answer to "does this solve the FM vocabulary-size problem
better than the alternatives listed in the Introduction."
\item What the $\lambda=0$ result (Section~\ref{sec:results:lambda}) implies
practically: since $\lambda$ trades off raw coverage against
proximity/conciseness, is there a principled way to pick $\lambda$ for a
given downstream FM, or is it genuinely a design knob left to the
practitioner (tying back to your Results outline's closing point)?
\item Revisit the controllability/interpretability claim from
Section~\ref{sec:introduction}'s contribution list, and the expert
acceptability results from Section~\ref{sec:method:expert-eval}: what does a
reader/expert actually get to inspect or judge that a codebook-based
tokenizer would not let them inspect? A short worked example (reusing the
context-tree figure from Section~\ref{sec:method:tokenizer}) would make this
concrete rather than asserted.
\end{itemize}
